\documentclass[11pt]{article}
\usepackage[margin=1in]{geometry}
\usepackage[T1]{fontenc}
\usepackage[utf8]{inputenc}
\usepackage{amsmath,amssymb,amsthm,mathtools}
\usepackage{microtype}
\usepackage[numbers,sort&compress]{natbib}
\usepackage[colorlinks=true,linkcolor=blue,citecolor=blue,urlcolor=blue]{hyperref}
\allowdisplaybreaks[2]
\emergencystretch=2em
\newtheorem{result}{Result}
\newcommand{\R}{\mathbb{R}}
\newcommand{\E}{\mathbb{E}}
\newcommand{\KL}{\operatorname{KL}}
\newcommand{\SPD}{\operatorname{SPD}}
\newcommand{\GL}{\operatorname{GL}}
\newcommand{\Cov}{\operatorname{Cov}}
\newcommand{\tr}{\operatorname{tr}}
\newcommand{\svec}{\operatorname{svec}}
\newcommand{\softmax}{\operatorname{softmax}}
\newcommand{\diag}{\operatorname{diag}}
\title{Gauge-Covariant Variational Attention:\\Independent Query, Key, Value, and Output Maps\\and the Limits of a Transformer Correspondence}
\author{Robert C. Dennis}
\date{September 4, 2026}
\begin{document}
\maketitle

\begin{abstract}
A transported-Gaussian variational free energy ties relevance scores and information messages to the same precision geometry. This restriction explains why expanding an inverse-transpose transport does not produce arbitrary independent query, key, value, and output projections. We characterize the original square-factor restriction and construct an explicit extension in which categorical source inference determines relevance while a separate Gaussian conditional predicts content. Every finite attention head with arbitrary rectangular, possibly rank-deficient projections is exactly the posterior predictive mean of this normalized conditional model. A structured variational family recovers its source posterior, and learned fiber morphisms provide a passive gauge-covariant formulation with explicit density and reference-measure bookkeeping. Independently, flat full-Gaussian attention admits an exact dot-product representation in Gaussian sufficient statistics, reducing a direct pairwise-factorization arithmetic bound without discarding covariance terms. Further results distinguish precision pooling from mixture uncertainty, show that unrestricted frames can encode full covariance even with identity local covariance, and give a same-state natural-gradient obstruction for arbitrary attention increments. These results establish an attention-level correspondence through additional probabilistic structure. They do not derive arbitrary projections, nonlinear transformer composition, or next-token cross entropy from the original single-channel free energy. No measured speedup, perplexity improvement, or priority claim is made.
\end{abstract}

\section{The correspondence to be established}
\label{sec:introduction}
Scaled dot-product attention assigns a relevance weight to each available source, forms a weighted content average, and maps that average into the residual representation \citep{Vaswani2017}. A Gaussian interaction model can also produce an attention row by exponentiating negative relative entropy. The similarity leaves a substantive question: does the Gaussian construction contain the independent learned maps of ordinary attention, or only a constrained subclass in which comparison and transmitted information share the same geometry?

The distinction concerns the inference problem. If several Gaussian sources describe one state, minimizing a weighted sum of divergences pools their natural parameters. The precision that determines compatibility also determines the information transmitted to the receiver. If instead a categorical variable selects which source supplies a prediction, the relevance posterior and the predicted content can be parameterized separately. The second interpretation permits independent value and output maps because source identity mediates a conditional content model. It introduces structure beyond the first interpretation.

This note develops both statements without equating them. Result~\ref{res:C1} characterizes the original linear factors. Result~\ref{res:C2} expresses the complete flat Gaussian score as a dot product of constrained sufficient statistics. Results~\ref{res:C3} and~\ref{res:C4} construct arbitrary four-map attention as normalized source inference followed by a gauge-covariant predictive readout. Results~\ref{res:C5} and~\ref{res:C6} delimit the dynamical and training-objective conclusions. Results~\ref{res:C7}--\ref{res:C9} establish the associated precision, uncertainty, and covariance-representation identities. The labels C1--C9 retain the correspondence to the accepted derivation package from which this manuscript was prepared.

Probabilistic accounts of attention already appear in the literature. Gaussian attention keys provide one relevant comparison \citep{Nguyen2022}, and attention as marginalization provides another \citep{SinghBuckley2023}. The variational identity used below is the usual evidence-minus-posterior-KL identity \citep{Blei2017}. The purpose here is to state the precise construction, all its normalizers, and its interface with transported Gaussian beliefs. The proofs do not assume that a probabilistic interpretation uniquely explains why a given architecture or loss should be selected.

\section{Setup and the restriction of one Gaussian channel}
\label{sec:setup}
All vectors are columns. Fix positive integers $n,d,k,v$, a finite residual context $X=(x_1,\ldots,x_n)$ with $x_j\in\R^d$, and arbitrary real matrices
\begin{equation}
W_Q,W_K\in\R^{k\times d},
W_V\in\R^{v\times d}, W_O\in\R^{d\times v}.
\label{eq:shapes}
\end{equation}
No rank assumption is imposed. For each receiver $i$, let $\mathcal J_i\subseteq\{1,\ldots,i\}$ be a nonempty finite causal support. A strictly causal first row requires a separately supplied initial token. Disallowed sources are absent from $\mathcal J_i$ rather than represented by infinite scores. Let $b_{ij}\in\R$ be a finite prescribed bias on this support and let $\tau>0$. The standard scaling is $\tau=\sqrt{k}$, but the arguments retain a general positive temperature. Define
\begin{align}
q_i&=W_Qx_i,& k_j&=W_Kx_j,& v_j&=W_Vx_j,\\
s_{ij}&=q_i^Tk_j/\tau+b_{ij},&
\mathcal Z_i(q_i)&=\sum_{j\in\mathcal J_i}e^{s_{ij}},&
a_{ij}&=e^{s_{ij}}/\mathcal Z_i(q_i).
\label{eq:head-definitions}
\end{align}
The head output is
\begin{equation}
y_i=W_O\sum_{j\in\mathcal J_i}a_{ij}v_j.
\label{eq:head-output}
\end{equation}
All quantities are finite in real arithmetic and $\mathcal Z_i(q_i)>0$. Numerical overflow or an implementation's clipping rules are separate questions.

For the original Gaussian model, write $\mathcal N(\mu_i,\Sigma_i)$ for a belief in local $p$-dimensional coordinates, with $\Sigma_i\in\SPD(p)$ and frame $U_i\in\GL(p)$. Its flat transport is
\begin{equation}
\Omega_{ij}=U_iU_j^{-1}.
\label{eq:flat-transport}
\end{equation}
The transported source has mean $\Omega_{ij}\mu_j$ and covariance $\Omega_{ij}\Sigma_j\Omega_{ij}^T$. The mean cross term in its comparison with the receiver is
\begin{equation}
\mu_i^T\Omega_{ij}^{-T}\Sigma_j^{-1}\mu_j
=(U_i^{-1}\mu_i)^T(U_j^T\Sigma_j^{-1}\mu_j).
\label{eq:old-cross-term}
\end{equation}
It resembles a query-key pairing but its factors obey an algebraic constraint.

\begin{result}[C1: exact characterization of the original factors]
\label{res:C1}
For square real matrices $A,B\in\R^{p\times p}$, there exist $U\in\GL(p)$ and $\Sigma\in\SPD(p)$ satisfying
\begin{equation}
A=U^{-1}, B=U^T\Sigma^{-1}
\label{eq:old-factors}
\end{equation}
if and only if $A$ is invertible and $A^TB\in\SPD(p)$.
\end{result}
\begin{proof}
If \eqref{eq:old-factors} holds, then $A$ is invertible and
\begin{equation}
A^TB=U^{-T}U^T\Sigma^{-1}=\Sigma^{-1}\in\SPD(p).
\end{equation}
Conversely, choose $U=A^{-1}$ and $\Sigma=(A^TB)^{-1}$. The first identity follows immediately, while
\begin{equation}
U^T\Sigma^{-1}=A^{-T}A^TB=B.
\end{equation}
The covariance is SPD by assumption, completing the construction.
\end{proof}

The result excludes singular and rectangular factors from this representation. It also excludes some invertible square pairs: $A=I$ and $B=\diag(1,-1)$ fail positivity. Different formulas for the two factors therefore do not confer arbitrary independent projection freedom. This is a restriction of the stated factorization, rather than a classification of every parameterization that computes the same attention function.

Transport imposes another restriction. Equation~\eqref{eq:flat-transport} gives
\begin{equation}
\Omega_{ii}=I,
\Omega_{ij}\Omega_{j\ell}=\Omega_{i\ell},
\Omega_{ji}=\Omega_{ij}^{-1}.
\label{eq:cocycle}
\end{equation}
A constant invertible transport on a self edge must be $I$. If the same transport occurs on all three edges of a transitive attended triple, its composition law gives $\Omega^2=\Omega$, hence again $\Omega=I$. These statements require the indicated edges; they do not apply to an isolated distinct-token edge with no self or composition constraint. Similarly, if $\Omega_{ij}^{-T}\Sigma_j^{-1}$ equals one constant matrix on every pair including self pairs, that matrix is SPD because its self-pair value is $\Sigma_i^{-1}$.

Neither restriction is removed by retaining only the cross term in \eqref{eq:old-cross-term}. The complete Gaussian divergence also contains source-dependent quadratic and covariance terms. The next section preserves them.

\section{The exact full-Gaussian sufficient-statistic representation}
\label{sec:lift}
Move every belief to reference coordinates by defining
\begin{equation}
\xi_i=U_i^{-1}\mu_i,
C_i=U_i^{-1}\Sigma_iU_i^{-T},
H_i=C_i^{-1}, S_i=C_i+\xi_i\xi_i^T.
\label{eq:common-frame}
\end{equation}
Here $S_i$ is the second moment matrix; $\mathcal J_i$ denotes source support. Invertible pushforward preserves relative entropy: density Jacobians cancel in the ratio, and the integration measure transforms with the inverse factor. Applying $U_i^{-1}$ to both receiver and transported source therefore gives
\begin{equation}
D_{ij}=\KL\left(\mathcal N(\xi_i,C_i)\middle\|\mathcal N(\xi_j,C_j)\right).
\label{eq:common-kl}
\end{equation}

\begin{result}[C2: sufficient-statistic attention and arithmetic bound]
\label{res:C2}
Under flat transport and exact pushforward, define $\svec$ on symmetric $p\times p$ matrices by retaining diagonal entries and multiplying each strictly upper-triangular entry by $\sqrt2$. Thus
\begin{equation}
\svec(A)^T\svec(B)=\tr(AB)
\label{eq:svec}
\end{equation}
for symmetric $A,B$. With lifted column features
\begin{equation}
\mathcal Q_i=\begin{pmatrix}\xi_i\\\svec(S_i)\\1\end{pmatrix},
\mathcal K_j=\begin{pmatrix}H_j\xi_j\\-\tfrac12\svec(H_j)\\
-\tfrac12(\xi_j^TH_j\xi_j+\log\det C_j)\end{pmatrix},
\label{eq:lifted-features}
\end{equation}
negative Gaussian relative entropy satisfies
\begin{equation}
-D_{ij}=\mathcal Q_i^T\mathcal K_j+\tfrac12(\log\det C_i+p).
\label{eq:lifted-score}
\end{equation}
For positive normalized source weights $\pi_{ij}$ on $\mathcal J_i$, consequently,
\begin{equation}
\begin{split}
\beta_{ij}&=\softmax_{j\in\mathcal J_i}(-D_{ij}/\tau+\log\pi_{ij})\\
&=\softmax_{j\in\mathcal J_i}(\mathcal Q_i^T\mathcal K_j/\tau+\log\pi_{ij}).
\end{split}
\label{eq:lifted-attention}
\end{equation}
The lifted dimension is $p+p(p+1)/2+1$. Dense preprocessing and all pair contractions admit the exact-arithmetic cost $O(np^3+n^2p^2)$, with $O(np^2)$ feature storage in addition to attention storage.
\end{result}
\begin{proof}
The Gaussian formula is
\begin{equation}
D_{ij}=\frac12\left[\log\frac{\det C_j}{\det C_i}-p+
\tr(H_jC_i)+(\xi_i-\xi_j)^TH_j(\xi_i-\xi_j)\right].
\end{equation}
Expanding the mean term yields
\begin{align}
-D_{ij}={}&\xi_i^TH_j\xi_j-\frac12\tr(S_iH_j)
-\frac12(\xi_j^TH_j\xi_j+\log\det C_j)\nonumber\\
&+\frac12(\log\det C_i+p).
\label{eq:expanded-gaussian}
\end{align}
The symmetric-vectorization identity follows by writing the Frobenius inner product as diagonal products plus twice the strictly upper-triangular products. Substitution proves \eqref{eq:lifted-score}. Its last term depends only on the receiver and cancels in a row softmax.

Dense per-token changes of frame, covariance factorizations, inversions, and log determinants cost $O(p^3)$ per token under the usual arithmetic model. Each feature has $O(p^2)$ entries and each pair contraction costs $O(p^2)$. Summing these costs gives the bound and storage requirement.
\end{proof}

This can remove the pairwise factorization term of a direct $O(n^2p^3)$ full-Gaussian score computation. It does not remove pairwise attention and is not a runtime measurement. A dot-product kernel must retain $1/\tau$ rather than silently scale by the square root of the enlarged feature dimension. Score clipping, regularization, backward rules, and finite-precision conditioning require separate implementation checks before claiming equivalence of an optimized path.

The lifted features are nonlinear, constrained functions of a belief. They are not arbitrary learned linear projections. Their structure has an exponential-family explanation. For $p_\eta(z)=h(z)\exp(\eta^TT(z)-A(\eta))$, write $m_i=\E_{q_i}T(z)$ and assume the displayed expectations and relative entropy exist. Direct substitution gives
\begin{equation}
-\KL(q_i\|p_{\eta_j})=m_i^T\eta_j-A(\eta_j)+\mathcal H(q_i)+\E_{q_i}\log h(z),
\label{eq:exponential-kl}
\end{equation}
where $\mathcal H$ is differential entropy. An expectation parameter supplies the query, a natural parameter supplies the key, and a negative log partition supplies the source bias. These quantities are finite for the nonsingular Gaussians used here.

The common frame is essential. Arbitrary curved edge transports, receiver-dependent projections, and diagonalization after dense transport do not automatically admit this factorization. Locally diagonal covariances generally become dense in the common frame. A reduced feature dimension $2p+1$ is available when covariances are diagonal in the common frame itself. Section~\ref{sec:diagonal} distinguishes representation capacity from covariance-family closure.
\section{Separating source inference from content prediction}
\label{sec:probabilistic}
The independent maps in \eqref{eq:shapes} can be realized by introducing separate comparison and content variables. Condition throughout this section on the deterministic context $X$, support, biases, and parameters. The comparison variable is a hypothetical auxiliary observation evaluated at $q_i$. It is not asserted to be a separately generated raw token while the same raw token is already fixed in $X$.

For $R\in\SPD(m)$, write the normalized Gaussian density as
\begin{equation}
\varphi_R(z)=(2\pi)^{-m/2}(\det R)^{-1/2}
\exp\left(-\frac12z^TR^{-1}z\right).
\label{eq:normal-density}
\end{equation}
Continuous densities initially use coordinate Lebesgue measure; a source label uses counting measure.

\begin{result}[C3: arbitrary four-map attention as a posterior mean]
\label{res:C3}
For every head satisfying Section~\ref{sec:setup}, choose $R\in\SPD(v)$ and define
\begin{equation}
D_i=\sum_{j\in\mathcal J_i}\exp\left(b_{ij}+\frac{\|k_j\|^2}{2\tau}\right),
\pi_{ij}=D_i^{-1}\exp\left(b_{ij}+\frac{\|k_j\|^2}{2\tau}\right).
\label{eq:compensating-prior}
\end{equation}
Introduce $J_i\in\mathcal J_i$, $Y_i\in\R^k$, and $Z_i\in\R^v$ with conditional joint density
\begin{equation}
p_i(j,y,z\mid X)=\pi_{ij}\varphi_{\tau I_k}(y-k_j)\varphi_R(z-v_j).
\label{eq:joint}
\end{equation}
This model is normalized. At $Y_i=q_i$, its exact source posterior is $a_i$, and
\begin{equation}
W_O\E[Z_i\mid Y_i=q_i,X]=y_i.
\label{eq:posterior-mean-head}
\end{equation}
For the structured variational family
\begin{equation}
q_r(j,z)=r_j\varphi_R(z-v_j), r_j\geq0,\sum_jr_j=1,
\label{eq:structured-family}
\end{equation}
the complete negative evidence lower bound is
\begin{equation}
\mathcal F_i(r;q_i)=\sum_jr_j\log\frac{r_j}{\pi_{ij}\varphi_{\tau I_k}(q_i-k_j)}
=\KL(r\|a_i)-\log p_i(q_i\mid X).
\label{eq:structured-elbo}
\end{equation}
Its unique minimizing row is $r=a_i$.
\end{result}
\begin{proof}
Every Gaussian factor integrates to one and the finite prior sums to one, proving normalization. Expansion of the comparison energy gives the pointwise identity
\begin{equation}
\pi_{ij}\varphi_{\tau I_k}(y-k_j)
=\frac{(2\pi\tau)^{-k/2}e^{-\|y\|^2/(2\tau)}}{D_i}
\exp(y^Tk_j/\tau+b_{ij}).
\label{eq:normalizer-cancellation}
\end{equation}
Define $\mathcal Z_i(y)=\sum_j\exp(y^Tk_j/\tau+b_{ij})$. The strictly positive evidence density is
\begin{equation}
p_i(y\mid X)=\frac{(2\pi\tau)^{-k/2}e^{-\|y\|^2/(2\tau)}\mathcal Z_i(y)}{D_i}.
\label{eq:query-evidence}
\end{equation}
Bayes' rule using these continuous density representatives gives
\begin{equation}
p_i(j,z\mid Y_i=q_i,X)=a_{ij}\varphi_R(z-v_j).
\label{eq:full-posterior}
\end{equation}
Its first content moment proves \eqref{eq:posterior-mean-head}.

The negative ELBO is $\E_{q_r}[\log q_r-\log p_i(J_i,q_i,Z_i\mid X)]$. The matched content log densities cancel, leaving the first expression in \eqref{eq:structured-elbo}, with $0\log0=0$. Substituting Bayes' formula gives the second expression. Categorical KL nonnegativity and its equality condition establish uniqueness.
\end{proof}

Every term can also be displayed without posterior notation:
\begin{align}
\mathcal F_i(r;q_i)={}&\frac{k}{2}\log(2\pi\tau)+\frac{\|q_i\|^2}{2\tau}+\log D_i\nonumber\\
&+\sum_jr_j\log r_j-\sum_jr_j(q_i^Tk_j/\tau+b_{ij}).
\label{eq:expanded-elbo}
\end{align}
The covariance is defined directly in content space. A rectangular or rank-deficient $W_V$ therefore creates no singular conditional density. Applying $W_O$ to content may produce a singular output pushforward, but it remains a valid probability measure with the stated mean; no nonsingular output density is asserted.

The prior \eqref{eq:compensating-prior} is a modeling addition. A uniform prior would give logits $q_i^Tk_j/\tau-\|k_j\|^2/(2\tau)$ up to receiver constants. For $q_i=0$, $k_1=0$, $k_2=2$, and $\tau=1$, dot-product attention assigns equal weights while uniform-prior Gaussian routing does not. Equal key norms remove that discrepancy, but arbitrary learned projections do not guarantee them. The declared prior cancels the source energy explicitly, and its normalizer remains in \eqref{eq:query-evidence}.

More generally, if $p(y\mid j)=h(y)\exp(T(y)^T\eta_j-A(\eta_j))$ is a normalized exponential-family likelihood with finite $A(\eta_j)$, choosing $\pi_j\propto\exp(A(\eta_j)+b_j)$ gives
\begin{equation}
p(j\mid y)=\softmax_j(T(y)^T\eta_j+b_j)
\label{eq:exponential-routing}
\end{equation}
wherever evidence is positive. The query supplies an observed statistic, the key a component parameter, the value conditional content, and the output map the space in which that content is used. These roles are compatible with variational inference but are not selected by the original one-belief construction.

\subsection{Why the variational family and readout matter}
\label{sec:meanfield}
Replacing the structured family by $r_jq(z)$ changes the variational problem. For $q(z)=\mathcal N(m,C)$, the exact gap to log evidence is
\begin{equation}
\KL(r\|a_i)+\sum_jr_j\KL\left(\mathcal N(m,C)\middle\|\mathcal N(v_j,R)\right).
\label{eq:mf-gap}
\end{equation}
At fixed $r$, Gaussian KL minimization gives $C=R$ and $m=\sum_jr_jv_j$. At fixed $m,C$, the source coordinate instead obeys
\begin{equation}
r_j\propto a_{ij}\exp\left[-\frac12(m-v_j)^TR^{-1}(m-v_j)\right].
\label{eq:mf-routing}
\end{equation}
The covariance trace term is common to the row; the content penalty generally is not. After optimizing the content factor, the extra cost is
\begin{equation}
\frac12\sum_jr_j(v_j-\bar v_r)^TR^{-1}(v_j-\bar v_r),
\bar v_r=\sum_jr_jv_j.
\end{equation}
These identities follow by substitution of the Gaussian KL formula into \eqref{eq:mf-gap}. Full product-family equality with the posterior requires all positive-weight content conditionals to coincide. Equality of one coordinate update with an attention mean is weaker.

A predictive mean also need not be a maximum a posteriori estimate. The equal mixture of $\mathcal N(-2,1)$ and $\mathcal N(2,1)$ has mean zero, but its density has second derivative $3\varphi_1(2)>0$ there. Its mean is a local density minimum. Result~\ref{res:C3} concerns a posterior moment.

\section{The moving-frame and probability-measure interface}
\label{sec:gauge}
At token $i$, introduce a residual fiber $E_i$ of dimension $d$, a comparison fiber $C_i$ of dimension $k$, and a content fiber $F_i$ of dimension $v$. These can be associated to the defining representations of $\GL(d)\times\GL(k)\times\GL(v)$, with the comparison dual representation also present. A finite token base admits a global trivialization. Edge transports are additional data; the graph alone does not specify a connection.

The learned maps and actual transports have types
\begin{equation}
Q_i:E_i\to C_i^*, K_i:E_i\to C_i,
V_i:E_i\to F_i, O_i:F_i\to E_i,
\label{eq:typed-maps}
\end{equation}
\begin{equation}
T^C_{ij}:C_j\to C_i, T^F_{ij}:F_j\to F_i.
\label{eq:typed-transports}
\end{equation}
Transports are invertible; learned maps may have any rank. A query is a covector, so its pairing with a transported key needs no unstated vector-dual identification.

\begin{result}[C4: covariant independent maps and a normalized Gaussian interface]
\label{res:C4}
Under independent passive coordinate changes $x_i'=g_i^Ex_i$, with invertible $g_i^E,g_i^C,g_i^F$ acting in their named fibers, let
\begin{align}
Q_i'&=(g_i^C)^{-T}Q_i(g_i^E)^{-1},& K_i'&=g_i^CK_i(g_i^E)^{-1},\\
V_i'&=g_i^FV_i(g_i^E)^{-1},& O_i'&=g_i^EO_i(g_i^F)^{-1},\\
(T^C_{ij})'&=g_i^CT^C_{ij}(g_j^C)^{-1},&
(T^F_{ij})'&=g_i^FT^F_{ij}(g_j^F)^{-1}.
\label{eq:gauge-laws}
\end{align}
For invariant scalar biases, the scores and output
\begin{equation}
s_{ij}=(Q_ix_i)^TT^C_{ij}K_jx_j/\tau+b_{ij},
\Delta x_i=O_i\sum_ja_{ij}T^F_{ij}V_jx_j
\label{eq:covariant-head}
\end{equation}
obey $s_{ij}'=s_{ij}$ and $\Delta x_i'=g_i^E\Delta x_i$. A flat synchronized choice realizes every head in \eqref{eq:head-output}. With transforming SPD comparison and content covariances, this construction admits the normalized source/content model of Result~\ref{res:C3}. Its probability law and categorical posterior are covariant; an absolute negative log evidence is an invariant scalar only after an invariant reference measure is specified.
\end{result}
\begin{proof}
Substituting the transformation laws into a score cancels the comparison changes at both endpoints. Thus scores and weights are invariant. A transported content vector transforms by $g_i^F$, which cancels its inverse inside $O_i'$, leaving $g_i^E\Delta x_i$. Residual addition consequently transforms in $E_i$.

Choose arbitrary reference matrices $\overline Q,\overline K,\overline V,\overline O$ and frames $U_i^E,U_i^C,U_i^F$. Set
\begin{align}
Q_i&=(U_i^C)^{-T}\overline Q(U_i^E)^{-1},&
K_i&=U_i^C\overline K(U_i^E)^{-1},\\
V_i&=U_i^F\overline V(U_i^E)^{-1},&
O_i&=U_i^E\overline O(U_i^F)^{-1},\\
T^C_{ij}&=U_i^C(U_j^C)^{-1},&T^F_{ij}&=U_i^F(U_j^F)^{-1}.
\label{eq:flat-witness}
\end{align}
With $\bar x_i=(U_i^E)^{-1}x_i$, cancellation gives
\begin{equation}
s_{ij}=(\overline Q\bar x_i)^T(\overline K\bar x_j)/\tau+b_{ij},
\Delta x_i=U_i^E\overline O\sum_ja_{ij}\overline V\bar x_j.
\end{equation}
Identity frames yield the standard head. No learned projection was inverted.

At receiver $i$, write
\begin{equation}
q_i=Q_ix_i, \zeta_i=H_i^{-1}q_i, t_{ij}=T^C_{ij}K_jx_j,
\end{equation}
where $q_i\in C_i^*$, $\zeta_i,t_{ij}\in C_i$, and $H_i\in\SPD(k)$ transforms as $H_i'=(g_i^C)^{-T}H_i(g_i^C)^{-1}$. This comparison metric is separately specified, not assumed to be a precision from \eqref{eq:common-frame}. Define
\begin{equation}
D_i^H=\sum_j\exp\left(b_{ij}+\frac{t_{ij}^TH_it_{ij}}{2\tau}\right),
\pi_{ij}^H=(D_i^H)^{-1}\exp\left(b_{ij}+\frac{t_{ij}^TH_it_{ij}}{2\tau}\right).
\end{equation}
The Lebesgue comparison likelihood is
\begin{equation}
p(y\mid j,X)=\frac{\sqrt{\det H_i}}{(2\pi\tau)^{k/2}}
\exp\left[-\frac{(y-t_{ij})^TH_i(y-t_{ij})}{2\tau}\right].
\label{eq:metric-gaussian}
\end{equation}
It is the normalized density of $\mathcal N(t_{ij},\tau H_i^{-1})$. At $y=\zeta_i$, expansion cancels the prior source energy and leaves \eqref{eq:covariant-head}. Its evidence is
\begin{equation}
p(\zeta_i\mid X)=\frac{\sqrt{\det H_i}}{(2\pi\tau)^{k/2}}
\exp\left[-\frac{q_i^TH_i^{-1}q_i}{2\tau}\right]
\frac{\sum_j e^{s_{ij}}}{D_i^H}.
\label{eq:metric-evidence}
\end{equation}
If the metric is parallel, $H_j=(T^C_{ij})^TH_iT^C_{ij}$, its source energy equals $(K_jx_j)^TH_j(K_jx_j)$. Using the transported energy directly does not require parallelism.

For content, set $m_{ij}=T^F_{ij}V_jx_j$, choose $G_i\in\SPD(v)$ transforming as a precision on $F_i$, and choose $\epsilon>0$. The conditional density
\begin{equation}
p(z\mid j,X)=\frac{\sqrt{\det G_i}}{(2\pi\epsilon)^{v/2}}
\exp\left[-\frac{(z-m_{ij})^TG_i(z-m_{ij})}{2\epsilon}\right]
\label{eq:metric-content}
\end{equation}
is normalized and independent of the comparison observation given $j,X$. The three factors define a proper joint. The structured content conditional cancels in its ELBO as before, and $O_i\E[Z_i\mid\zeta_i,X]$ equals the covariant head output.

Let $g=g_i^C$. Under $y'=gy$ and $H_i'=g^{-T}H_ig^{-1}$, the quadratic is invariant and
\begin{equation}
\sqrt{\det H_i'}=\sqrt{\det H_i}/|\det g|.
\end{equation}
Thus $p'(y'\mid j,X)=p(y\mid j,X)/|\det g|$, while $\mathrm{d}y'=|\det g|\mathrm{d}y$. This is exactly the change of density for the same probability law. The common Jacobian cancels from source posterior ratios. Content has the corresponding factor $|\det g_i^F|^{-1}$.

With Lebesgue densities, the negative ELBO transforms as
\begin{equation}
\mathcal F_i'(r;g\zeta_i)=\mathcal F_i(r;\zeta_i)+\log|\det g|.
\label{eq:elbo-jacobian}
\end{equation}
Its posterior KL gap remains invariant. Alternatively, declare metric volume
\begin{equation}
\mathrm{d}\operatorname{vol}_{H_i}(y)=\sqrt{\det H_i}\mathrm{d}y.
\label{eq:metric-volume}
\end{equation}
This measure is invariant and the Gaussian density with respect to it has normalization $(2\pi\tau)^{-k/2}$. The density and negative log evidence are scalars in that convention. The reference-measure qualification follows.
\end{proof}

The geometric freedom resides in learned morphisms as well as connections. The effective comparison $Q_i^TT^C_{ij}K_j$ has rank at most $k$, and $O_iT^F_{ij}V_j$ has rank at most $v$. Neither generally obeys a transport cocycle. When $v<d$, the latter cannot be the identity on a self edge. In a flat reference gauge,
\begin{equation}
\overline Q\mapsto A^{-T}\overline Q,
\overline K\mapsto A\overline K,
\overline V\mapsto B\overline V,
\overline O\mapsto\overline OB^{-1}
\end{equation}
leave the head function unchanged for appropriate invertible $A,B$. Four parameter factors do not imply four independently identifiable observable functions.

Noise is specified in its named fiber. If $S\in\SPD(d)$, then $WSW^T\in\SPD(k)$ exactly when $W\in\R^{k\times d}$ has full row rank: its quadratic form vanishes exactly on $\ker W^T$. Separately added SPD noise repairs singularity if it transforms with the coordinates. Even a full-row-rank rectangular projection need not preserve KL. For $W=(1,0)$, unit-covariance Gaussians in $\R^2$ with means $(0,1)^T$ and $(0,0)^T$ have KL $1/2$, while their projected laws coincide.

\subsection{Passive covariance, active symmetry, and shared weights}
\label{sec:active}
Result~\ref{res:C4} transforms every map and metric. Holding a numerical matrix $W$ fixed while inputs and outputs undergo an active action instead requires
\begin{equation}
W\rho_{\mathrm{in}}(g)=\rho_{\mathrm{out}}(g)W,\text{ for every }g.
\label{eq:intertwiner}
\end{equation}
For the defining $\GL(p)$ representation on both sides, commuting with all invertible diagonal matrices forces $W$ to be diagonal; commuting with elementary shears forces equal diagonal entries, so $W=cI$. For a map from the defining representation to its dual, $Wg=g^{-T}W$ and $g=2I$ force $W=0$. A transforming metric or different representation supplies additional structure; passive covariance does not evade these restrictions.

Connection-preserving weight sharing is stronger than covariance. A morphism $L_i:E_i\to F_i$ preserves transport if $L_iT^E_{ij}=T^F_{ij}L_j$. On a connected graph with inverse reverse transports, a basepoint map $L_0$ extends independently of path precisely when
\begin{equation}
L_0\mathcal H^E_\gamma=\mathcal H^F_\gamma L_0
\label{eq:holonomy}
\end{equation}
for every based loop $\gamma$, where $\mathcal H_\gamma$ is its holonomy. Necessity follows by composing edge conditions around a loop. For sufficiency, define $L_i=T^F_{i0}L_0(T^E_{i0})^{-1}$ along a path; comparing paths reduces to \eqref{eq:holonomy}. Trivial holonomy permits arbitrary rectangular $L_0$. The flat witness does not establish arbitrary connection-preserving weights for unspecified curved geometry.

For multiple heads, take comparison and content fibers for each head and apply the construction to each source posterior. A linear map from their direct sum to the residual fiber implements concatenation followed by output projection. This proves the attention and linear recombination correspondence; it supplies no transformation rule for an undeclared normalization, feedforward nonlinearity, positional operation, or decoder.
\section{A dynamical obstruction and a valid conditional update}
\label{sec:gradient}
An attention output can be a prediction without being a natural-gradient direction on the query state. For fixed keys and values of the same dimension as a variable query $q$, define
\begin{equation}
a_j(q)=\softmax_j(q^Tk_j/\tau), f(q)=\sum_ja_j(q)v_j.
\end{equation}
Differentiating the finite sum gives
\begin{equation}
Df(q)=\frac1\tau\left[\sum_ja_jv_jk_j^T-
\left(\sum_ja_jv_j\right)\left(\sum_ja_jk_j\right)^T\right].
\label{eq:attention-jacobian}
\end{equation}
Independent values need not make this Jacobian compatible with a scalar potential.

\begin{result}[C5: no universal same-state natural-gradient interpretation]
\label{res:C5}
There exists a two-source attention increment $f:\R^2\to\R^2$ for which no $C^2$ scalar $E$ and smooth SPD metric $G(q)$ satisfy $f(q)=-G(q)^{-1}\nabla E(q)$ on a neighborhood of zero.
\end{result}
\begin{proof}
Set $\tau=1$, use zero biases, and choose
\begin{equation}
k_1=\begin{pmatrix}1\\0\end{pmatrix}, k_2=\begin{pmatrix}-1\\0\end{pmatrix},
v_1=\begin{pmatrix}0\\1\end{pmatrix}, v_2=\begin{pmatrix}0\\-1\end{pmatrix}.
\end{equation}
Then $f(q)=(0,\tanh q_1)^T$, so $f(0)=0$ and
\begin{equation}
Df(0)=N=\begin{pmatrix}0&0\\1&0\end{pmatrix}.
\end{equation}
Suppose a representation existed. The identity $\nabla E=-Gf$ implies $\nabla E(0)=0$. Differentiating at zero removes every derivative-of-metric term because $f(0)=0$, leaving
\begin{equation}
\nabla^2E(0)=-G(0)N.
\end{equation}
Write $G(0)=\left(\begin{smallmatrix}a&b\\b&c\end{smallmatrix}\right)$. Then $G(0)N=\left(\begin{smallmatrix}b&0\\c&0\end{smallmatrix}\right)$. Hessian symmetry forces $c=0$, contradicting positive definiteness. This applies to state-dependent as well as constant metrics.
\end{proof}

The result concerns the increment in $q\mapsto q+\eta f(q)$ on the same query state. It does not exclude a larger latent state, conditional prediction, different damping, or an overwrite map whose increment would be $f(q)-q$. Tied special cases are possible: $\sum_ja_j(q)k_j$ is the gradient of $\tau\log\sum_j\exp(q^Tk_j/\tau)$.

A separate content state admits a natural-gradient update when routing is fixed. With $\bar v=\sum_ja_jv_j$ and $R\in\SPD(v)$, define
\begin{equation}
F(m)=\frac12\sum_ja_j(m-v_j)^TR^{-1}(m-v_j).
\end{equation}
The fixed-covariance Gaussian mean family has Fisher metric $R^{-1}$, so
\begin{equation}
\nabla_mF=R^{-1}(m-\bar v),
\operatorname{grad}_{Fisher}F=m-\bar v.
\end{equation}
A unit Euler step sends every $m$ to $\bar v$; a general step $\eta$ gives $(1-\eta)m+\eta\bar v$. This is a declared content-coordinate problem with frozen routing. Differentiating through changing weights or identifying $m$ with the query changes its premises.

\section{Auxiliary evidence and language-model training}
\label{sec:objective}
At the solution of Result~\ref{res:C3}, the complete negative log evidence is
\begin{align}
-\log p_i(q_i\mid X)={}&\frac{k}{2}\log(2\pi\tau)+\frac{\|q_i\|^2}{2\tau}
+\log\sum_j\exp\left(b_{ij}+\frac{\|k_j\|^2}{2\tau}\right)\nonumber\\
&-\log\sum_j\exp\left(b_{ij}+\frac{q_i^Tk_j}{\tau}\right).
\label{eq:optimized-evidence}
\end{align}
The source-prior normalizer and the attention normalizer both remain.

\begin{result}[C6: local comparison evidence is not generally cross entropy]
\label{res:C6}
At fixed context, \eqref{eq:optimized-evidence} is independent of $W_V,W_O$. It is not generally equal to next-token cross entropy, even up to a parameter-independent additive constant.
\end{result}
\begin{proof}
Integrating the normalized content kernel removes it from the comparison evidence. Equation~\eqref{eq:optimized-evidence} consequently contains no value or output map. For a counterexample to a general objective identity, take a one-source row with $q_i=k_1=0$, zero bias, and fixed temperature. Vary content/output parameters so a binary decoder has logits $(t,0)$, with the first class as target. Comparison evidence is constant in $t$, but cross entropy is $\log(1+e^{-t})$, which is not constant. The two objectives cannot agree up to a parameter-independent constant for all admitted heads.
\end{proof}

The fixed-context qualification is local. Earlier values can influence later queries, but that indirect dependence does not remove the counterexample or prove a whole-stack objective identity. Learning content maps for prediction requires an additional declared objective or observation model.

An informative target generally changes source inference. For target likelihood $p(t\mid z,X)$, Bayes' rule gives
\begin{equation}
p(J_i=j\mid q_i,t,X)\propto a_{ij}\int p(t\mid z,X)\varphi_R(z-v_j)\mathrm{d}z.
\label{eq:target-posterior}
\end{equation}
The extra factor depends on the source whenever predictive evidence differs. Forward attention may remain a target-blind predictive computation, but it is not thereby exact inference conditioned on that informative target.

Nonlinear computation also need not commute with replacing content by its mean. Let a routed scalar logit $T$ equal $0$ or $2$ with equal probabilities, and let $\sigma$ be logistic. Then
\begin{equation}
\sigma(\E T)=\frac{e}{1+e},
\E\sigma(T)=\frac12\left(\frac12+\frac{e^2}{1+e^2}\right),
\end{equation}
and the first minus the second is
\begin{equation}
\frac{(e-1)^3}{4(1+e)(1+e^2)}>0.
\label{eq:nonlinear-difference}
\end{equation}
Here $T$ can be the selected component's mean, so the example does not require a zero-covariance density. Attention-mean equality proves neither equality of marginalized predictions nor equality of training gradients through an arbitrary nonlinear stochastic hierarchy.

\subsection{An exact cross-entropy embedding with limited explanatory force}
\label{sec:manufactured}
These distinctions do not rule out every normalized ELBO representation of a transformer. A deliberately supplied construction shows why representability alone is too weak to complete the derivation.

Let $p_0(X)$ be a normalized positive density on finite continuous residual contexts, independent of trainable parameters. For each row let $m_i=|\mathcal J_i|$, choose $J_i$ uniformly, and introduce Bernoulli $B_i$ with
\begin{equation}
p_\theta(B_i=1\mid J_i=j,X)=a_{ij}(X),
p_\theta(B_i=0\mid J_i=j,X)=1-a_{ij}(X).
\label{eq:manufactured-bernoulli}
\end{equation}
Both outcomes belong to the model. Give content a normalized Gaussian kernel $\varphi_R(z_i-v_{J_i}(X))$. Conditional independence across rows gives $p(B_i=1\mid X)=1/m_i$, and conditioning on $B_i=1$ yields source posterior $a_i(X)$.

Separately supply a normalized finite-vocabulary conditional $d_\theta(t\mid X)$ with finite logits computed by a well-defined deterministic transformer. Normalizations, feedforward operations, residual paths, and attention means are already contained in this decoder. Declare the target independent of realized auxiliary sources and content given $X$. The full joint is
\begin{equation}
\begin{split}
p_\theta(X,t,J,B,Z)={}&p_0(X)d_\theta(t\mid X)\\
&\times\prod_i\left[\frac1{m_i}p_\theta(B_i\mid J_i,X)
\varphi_R(z_i-v_{J_i}(X))\right].
\end{split}
\label{eq:manufactured-joint}
\end{equation}
Every factor normalizes conditionally, hence the joint normalizes. At $B_i=1$ for all rows,
\begin{equation}
p_\theta(X,t,B=\mathbf1)=p_0(X)d_\theta(t\mid X)\prod_i m_i^{-1}.
\label{eq:manufactured-evidence}
\end{equation}
For a product of structured row families, cancellation gives the entire ELBO
\begin{equation}
\mathcal L=\log p_0(X)+\log d_\theta(t\mid X)-\sum_i\log m_i-\sum_i\KL(r_i\|a_i).
\label{eq:manufactured-elbo}
\end{equation}
At $r_i=a_i$, negative ELBO equals cross entropy plus $-\log p_0(X)+\sum_i\log m_i$, independent of parameters when masks and $p_0$ are fixed. The optimum value and its parameter derivative therefore match the supplied cross entropy up to this constant. The conditional construction can accompany a deterministic autoregressive computation whose categorical chain-rule product is already a normalized sequence model.

This embedding contains attention in the Bernoulli likelihood and the complete transformer in the target conditional. Auxiliary samples add no predictive information about $t$ given $X$. It establishes a weak representation and prevents a universal nonexistence claim for cross-entropy embeddings. It does not derive the decoder, projections, or loss from the original transported-Gaussian functional. Removing the supplied transformer decoder removes the asserted predictive law.

A full correspondence still requires one normalized model for the actual observed variables and latent states, its variational family, all partition terms, and a proof connecting inference and objective to the entire deterministic network. Every nonattention map also needs a domain and coordinate-transformation rule before global covariance can be claimed. Finite inputs do not guarantee that an unspecified normalization is defined.

\section{Consensus, source prediction, and uncertainty}
\label{sec:uncertainty}
The original Gaussian interaction has a precise interpretation different from source selection. It explains the inverse transpose without treating it as a missing product of independent weights.

\begin{result}[C7: weighted Gaussian KL minimization pools natural parameters]
\label{res:C7}
Let $p_j=\mathcal N(\xi_j,C_j)$ be fixed nonsingular Gaussian sources, $H_j=C_j^{-1}$, and $\beta_j\geq0$ fixed weights summing to one. Over probability densities for which the objective is finite, the unique minimizer of
\begin{equation}
\mathcal F(q)=\sum_j\beta_j\KL(q\|p_j)
\label{eq:pool-objective}
\end{equation}
is Gaussian with
\begin{equation}
H^*=\sum_j\beta_jH_j, h^*=\sum_j\beta_jH_j\xi_j,
\xi^*=(H^*)^{-1}h^*.
\label{eq:precision-pool}
\end{equation}
Sources and weights remain frozen in this coordinate problem.
\end{result}
\begin{proof}
Set $\widetilde p(z)=\prod_jp_j(z)^{\beta_j}$. Its logarithm is quadratic with precision $H^*$. For nonzero $u$, $u^TH^*u=\sum_j\beta_ju^TH_ju>0$, since some weight is positive. Thus $H^*$ is SPD and $Z=\int\widetilde p(z)\mathrm{d}z$ is finite and positive. Expansion gives
\begin{equation}
\mathcal F(q)=\KL\left(q\middle\|Z^{-1}\widetilde p\right)-\log Z.
\end{equation}
KL nonnegativity and its equality condition give the unique normalized minimizer. Completing the quadratic square yields \eqref{eq:precision-pool}.
\end{proof}

If every $H_j=H$, the mean becomes $\sum_j\beta_j\xi_j$. Otherwise a solve is required. In receiver coordinates, the source precision and information-vector contributions are
\begin{equation}
\Omega_{ij}^{-T}\Sigma_j^{-1}\Omega_{ij}^{-1},
\Omega_{ij}^{-T}\Sigma_j^{-1}\mu_j.
\label{eq:local-messages}
\end{equation}
These are the transported natural parameters of the same Gaussian whose KL determined compatibility. The inverse transpose is fixed by that law and supplies no independent value projection.

For a prior KL coefficient $\alpha>0$ and source coefficient $\lambda>0$, the same proof gives
\begin{equation}
H^*=\frac{\alpha H_0+\lambda\sum_j\beta_jH_j}{\alpha+\lambda},
\end{equation}
\begin{equation}
\xi^*=\left(\alpha H_0+\lambda\sum_j\beta_jH_j\right)^{-1}
\left(\alpha H_0\xi_0+\lambda\sum_j\beta_jH_j\xi_j\right).
\label{eq:prior-pool}
\end{equation}
The precision denominator follows from the entropy coefficient of a weighted sum of KLs. A product likelihood with one entropy term defines a different objective. These fixed-source identities do not alone normalize an interacting population with shared random node states.

\begin{result}[C8: source-mixture predictive covariance]
\label{res:C8}
If $p(J=j\mid q)=a_j$ and $Z\mid J=j,q\sim\mathcal N(v_j,R_j)$ with $R_j\in\SPD(v)$, then
\begin{equation}
\bar v=\E[Z\mid q]=\sum_ja_jv_j,
\end{equation}
\begin{equation}
\Cov(Z\mid q)=\sum_ja_j\left[R_j+(v_j-\bar v)(v_j-\bar v)^T\right].
\label{eq:mixture-covariance}
\end{equation}
\end{result}
\begin{proof}
Conditioning on $J$ gives $\E[ZZ^T\mid q]=\sum_ja_j(R_j+v_jv_j^T)$. Subtracting $\bar v\bar v^T$ and using $\sum_ja_jv_j=\bar v$ yields the formula. Every between-source outer product is positive semidefinite.
\end{proof}

The between-source term represents uncertainty about which prediction applies. It may remain nonzero even when every conditional covariance is small. Squared weights would instead arise in a weighted sum of independent draws, a different random variable. Precision pooling estimates a state consistent with weighted sources, whereas source selection predicts content supplied by one source. Coincident means in special cases do not identify their uncertainty semantics.

\section{Diagonal covariance, coordinate closure, and representation}
\label{sec:diagonal}
Diagonal covariance can describe two different restrictions: closure of a family under coordinate changes, or the common-frame distributions expressible by combining a covariance with a learned frame. These properties are not equivalent.

An invertible $G$ sends every positive diagonal $D$ to a diagonal $GDG^T$ if and only if $G$ is monomial, a permutation times an invertible diagonal matrix. For necessity, each off-diagonal entry is $\sum_aG_{ra}G_{sa}D_{aa}$. Its vanishing for every positive diagonal choice forces $G_{ra}G_{sa}=0$ for every $a$ and $r\ne s$. Each column has at most one nonzero entry; invertibility forces exactly one per row and column. The converse follows directly. Generic general-linear congruences therefore do not preserve diagonal covariance in fixed coordinates.

\begin{result}[C9: full covariance represented by identity local covariance]
\label{res:C9}
For every $\xi\in\R^p$ and $C\in\SPD(p)$, an unrestricted full frame $U\in\GL(p)$ and local mean $\mu$ with local covariance $I$ represent the exact common-frame Gaussian $\mathcal N(\xi,C)$. The frame can be SPD and the exponential of a real symmetric matrix.
\end{result}
\begin{proof}
Choose the positive inverse square root $U=C^{-1/2}$ and $\mu=U\xi$. Then
\begin{equation}
U^{-1}\mu=\xi, U^{-1}IU^{-T}=C.
\end{equation}
This is the exact pushforward. The SPD matrix $U$ has a real symmetric logarithm by spectral decomposition, so $U=\exp(\log U)$.
\end{proof}

The result requires unrestricted per-token full frames and exact covariance pushforward. It establishes neither equal optimization dynamics nor equal parameter costs or predictive behavior for particular implementations. Restricted or shared frames, diagonal projection after transport, gradient filtering, decoder dependence on local coordinates, and computational budgets can change the comparison. Projection of a dense transported covariance onto its diagonal changes the distribution.

A perplexity difference between configurations labeled full and diagonal cannot alone establish an intrinsic representational ceiling of local diagonal covariance. Identifying its cause requires family definitions, frame restrictions, training objectives, evaluation protocols, and matched computational evidence. The algebra here does not supply that empirical identification.

\section{Discussion and the remaining derivation}
\label{sec:discussion}
The four transformer maps become distinct when relevance and content receive distinct probabilistic roles. A comparison statistic and component parameter determine a categorical posterior. A conditional content mean and residual readout determine its prediction. Learned morphisms permit these roles to coexist with coordinate covariance while genuine transports retain invertibility, duality, and composition. This is an exact construction for arbitrary finite heads, including singular and rectangular projections.

The original Gaussian channel retains its own exact attention correspondence. Its query contains expectation statistics and its key contains natural parameters and a log-partition term. This yields an uncertainty-preserving score representation and an arithmetic route for avoiding pairwise factorizations. It also explains the constrained factors: one distribution supplies both comparison geometry and information pooling. An independent content channel is an additional modeling choice rather than an algebraic rearrangement of that precision.

A whole-language-model theorem must identify observed tokens, auxiliary observations, latent sources and states, the variational family, and all likelihood and partition terms of one model. It must establish the relationship between inference or predictive moments and the complete deterministic network, including normalization, feedforward operations, residual composition, and the categorical decoder. If training equivalence is intended, equality of objectives and relevant derivatives must follow. An embedding that assumes the completed decoder, such as Section~\ref{sec:manufactured}, does not meet that stronger explanatory obligation.

The results suggest empirical questions rather than answering them. Implementing Result~\ref{res:C2} could test whether exact full-covariance scores become practical at larger width. A controlled comparison could distinguish precision pooling from independent content prediction at fixed routing. Uncertainty evaluation could compare Result~\ref{res:C8} with Result~\ref{res:C7}. None of these experiments was performed for this note. Representability alone implies no improvement in runtime, perplexity, calibration, or optimization.

\appendix
\section{Source package and scope of this write-up}
\label{app:provenance}
This note converts the September 4, 2026 derivation package in the V3\_Transformer repository under \texttt{docs/derivations/2026-09-04-gauge-vfe-transformer/}. The direct mathematical sources are \texttt{root-derivations.md}, \texttt{agent-gauge.md}, and \texttt{agent-probabilistic.md}. Accepted statements and limitations are recorded in \texttt{claim-ledger.json} and \texttt{theory-adjudication.md}. Results~\ref{res:C1}--\ref{res:C9} correspond to C1--C9 there. Their assumptions and proofs are reproduced here so the manuscript is self-contained.

The source package distinguishes accepted local mathematical results from an unresolved whole-language-model target. This write-up preserves that distinction and reports no new model tests, training runs, GPU benchmarks, or independent experimental verification. Saved run summaries motivated the question but do not establish a measured benefit of the extension. No perplexity endpoint is used as a proof of a covariance-family capacity limit. The cited literature supplies context for the displayed derivations, rather than substituting for their proofs; no systematic priority claim is made.

\begin{thebibliography}{9}
\bibitem[Vaswani et al.(2017)]{Vaswani2017}
A. Vaswani, N. Shazeer, N. Parmar, J. Uszkoreit, L. Jones, A. N. Gomez,
\L{}. Kaiser, and I. Polosukhin.
\newblock Attention Is All You Need.
\newblock \emph{Advances in Neural Information Processing Systems}, 30, 2017.
\newblock \url{https://arxiv.org/abs/1706.03762}.

\bibitem[Nguyen et al.(2022)]{Nguyen2022}
Tam Minh Nguyen, Tan Minh Nguyen, Dung D. D. Le, Duy Khuong Nguyen,
Viet-Anh Tran, Richard Baraniuk, Nhat Ho, and Stanley Osher.
\newblock Improving Transformers with Probabilistic Attention Keys.
\newblock \emph{Proceedings of the 39th International Conference on Machine Learning},
PMLR 162:16595--16621, 2022.
\newblock \url{https://proceedings.mlr.press/v162/nguyen22c.html}.

\bibitem[Singh and Buckley(2023)]{SinghBuckley2023}
Ryan Singh and Christopher L. Buckley.
\newblock Attention: Marginal Probability is All You Need?
\newblock arXiv:2304.04556, 2023.
\newblock \href{https://doi.org/10.48550/arXiv.2304.04556}{doi:10.48550/arXiv.2304.04556}.

\bibitem[Blei et al.(2017)]{Blei2017}
David M. Blei, Alp Kucukelbir, and Jon D. McAuliffe.
\newblock Variational inference: A review for statisticians.
\newblock \emph{Journal of the American Statistical Association}, 112(518):859--877, 2017.
\newblock \href{https://doi.org/10.1080/01621459.2017.1285773}{doi:10.1080/01621459.2017.1285773}.
\end{thebibliography}
\end{document}
